% Section 3: Results

\subsection{Introduction}
To be written.

\subsection{Performance on benchmark data set}
Paragraph on LLM identifying collusive transcripts from the benchmark data set.

\subsection{Human Audit}
The remainder of the section will summarize the findings from a human audit of the LLM’s output. The human auditor examined the 301 transcripts with the highest validated scores assigned by the LLM, which makes up .06\% of the corpus. With the LLM having determined these transcripts contain collusive content, the human auditor will make their own independent evaluation. If the human auditor finds the transcript contains collusive content then it is classified as a ``true positive'', and if the human auditor does not find the transcript contains collusive content then it is classified as a ``false positive''. 

Of course, the human auditor may be mistaken in their judgment – in the sense that the content did not result in firms coordinating their conduct to reduce competition -- but the goal is not to identify actual episodes of collusion but to determine if an LLM can replicate a human auditor.

The ensuing analysis addresses three questions. First, how accurate is an LLM in identifying collusive content? Accuracy will be measured by the true positive rate. This exercise will inform us of the usefulness of an LLM at substituting for human auditors in evaluating public communications. Second, what describes the collusive content of true positives? That is, what is the basis for concluding a firm’s announcement facilitates collusion? This analysis will advance our understanding of firm conduct and contribute to the body of work associated with Aryal et al (2022) and Harrington (2022). Third, when the LLM gets it wrong, how does it get it wrong? That will inform us of the LLM’s limitations and may suggest how to improve its performance.

\subsubsection{Human Audit Procedure}
Balancing the value of reviewing more transcripts against the time it takes to review them, we decided to review the top 300 transcripts – which makes up about one-half of the LLM-validated transcripts -- subject to including all transcripts with the same score. This resulted in 301 transcripts which comprise all those with a score of at least 76.92. The scores ranged from 76.92 to 85.00. The scores are summarized in Table X1 and Figure Y1.

% Placeholder for Table X1 and Figure Y1
\begin{center}
    [Insert Table X1 and Figure Y1 here]
\end{center}

For each of these transcripts, the human auditor read the excerpts provided by the LLM. In order to make for a more independent judgment, the human auditor did not read the LLM’s explanation for its score. The human auditor was aware that the LLM had given the transcript a high score and thus knew the LLM had rated the content as collusive. It is possible this knowledge could bias their evaluation though we have no reason to think it did.

The human auditor evaluated each transcript’s excerpts in terms of whether it would facilitate collusive effect. They did not evaluate it in terms of collusive intent. Intent is difficult to determine and our more general concern is whether public communications might reduce competition, whether intended or not. The LLM was instructed to consider intent, as well as effect, in order to better communicate what it is we are looking for. It is possible the LLM could have assigned a lower score where it finds evidence of effect but not intent; that is, the LLM assesses the content could reduce competition but concludes that was not the firm’s intent. However, based on our review of the LLM’s output, we do not believe its assessment was so nuanced.

Upon reading the excerpts, the human auditor concluded whether the content could facilitate firms jointly reducing competition such as by reducing supply or capacity or raising price. If that was determined to be the case then it was classified as a ``true positive'' (and labelled T), and otherwise as a ``false positive'' (and labelled as F). If T then the human auditor summarized the collusive content. If F, the human auditor summarized the content that might have led the LLM to give the transcript a high score. In some cases, there is no relevant content in which case it is denoted ``no content''. The LLM’s output along with the human auditor’s evaluation and summary are provided in Online Appendix A.

\subsection{LLM’s Performance}
Of the 301 transcripts with the highest LLM-validated score, the human auditor rated 109 as containing collusive content and 192 as lacking collusive content for a true positive rate of 36.2\%. Table X2 reports the true positive rate after partitioning the sample into three roughly equal-sized subsamples (while respecting score breaks). The true positive rate is decreasing in the score with a true positive rate of 41.5\% for the 106 transcripts in the highest bin and 31.9\% for the 94 transcripts in the lowest bin.

With details provided later, 163 of the 192 false positive transcripts contained content that, while relevant, did not cross the line for facilitating coordinated conduct. The other 29 false positive transcripts lacked any relevant content (``no content'') in which case it may be appropriate to classify them as hallucinations. With that definition, the hallucination rate is 9.6\% (29 out of 301).

\begin{table}[htbp]
    \centering
    \caption{LLM Performance Statistics}
    \begin{tabular}{l c c c c c}
        \hline
        Sample & Scores & Number & True positive & False positive & True positive rate \\
        \hline
        Full sample & 76.92-85.00 & 301 & 109 & 192 & 36.2\% \\
        Top bin & 79.58-85.00 & 106 & 44 & 62 & 41.5\% \\
        Middle bin & 77.92-79.58 & 101 & 35 & 66 & 34.7\% \\
        Bottom bin & 76.92-77.85 & 94 & 30 & 64 & 31.9\% \\
        \hline
    \end{tabular}
    \label{tab:X2}
\end{table}

Listed in Appendix A, there are 71 distinct companies in these 109 true positive transcripts and the four companies with the largest number of transcripts are Gol Linhas Aéreas Inteligentes (7 transcripts) which is a domestic airline in Brazil, Micron Technology (6) which is a manufacturer of computer memory and data storage, Norske Skogindustrier (4) which is a pulp and paper manufacturer, and Coca-Cola Bottlers Japan (4). We will later take a closer look at Gol Linhas Aéreas Inteligentes.

22 of the 109 transcripts are from previously documented episodes of companies making public announcements with collusive content, including 15 associated with the airlines case. 10 of those 15 airlines transcripts were from companies in the Aryal et al (2022) study: US Airways (3 transcripts), American Airlines (2), Delta Airlines (2), United Airlines (2), and AirTran (1); and then there were five transcripts from Air France-KLM (3) and Frontier Airlines (2). The other previously documented episodes are: auto rental (Avis), broiler chicken (Pilgrim’s Pride), generic drugs (Teva Pharmaceuticals), and pork (Smithfield Foods, WH Group).

In some of these cases – generic drugs for sure and possibly broiler chicken and pork – collusion also involved the more standard private communications. The LLM flagging them is a reminder that screening public communications for collusive content could lead to investigations that uncover the most compelling evidence of illegal collusion: direct, express, and private communications between competitors.

\subsection{Evaluation of Collusive Content}
Next we turn to investigating the content of firms’ communications. For the 109 true positive transcripts, Appendix B categorizes them according to the type of collusive content. Here, we focus on the five most common categories which are:
\begin{itemize}
    \item A firm announces the industry has reduced supply or capacity and references competitors or the industry in a manner supportive of such conduct (40 transcripts)
    \item A firm announces the industry needs to reduce supply or capacity or to show discipline (30 transcripts)
    \item A firm announces the industry has shown or will show supply or capacity or price discipline (12 transcripts)
    \item A firm announces the industry needs to raise prices (12 transcripts)
    \item A firm announces it is being or will be less aggressive in pricing and references competitors’ conduct in a manner encouraging similar conduct (8 transcripts)
\end{itemize}

\subsubsection{Category 1: The industry has reduced supply or capacity and references competitors or the industry in a manner supportive of such conduct}
The most common collusive content is that the industry has reduced supply or capacity and refers to competitors or the industry to facilitate such conduct. This type of content is present in 37\% of the 109 true positive transcripts. In emphasizing supply or capacity cuts, a firm may highlight how those reductions are in the industry’s best interests. For example:

\begin{quote}
``We've seen a massive … consolidation in the industry ... and hopefully, that leads to a better, a healthier competitive environment as well. No one is going to build DRAM capacity for the foreseeable future. It's just not going to happen.'' \\
$[508027]$
\end{quote}

\begin{quote}
``Other companies in our industry have announced considerable domestic downtime and capacity closures [and] share my philosophy that if you can't sell it or you can't use it, you shouldn't make it.'' \\
$[597966]$
\end{quote}

\begin{quote}
``We will not continue to invest at this pace, and our activity will slow. It's kind of a no-brainer from an economic standpoint. Our partners are making the right decisions in terms of the decreasing capital in this environment.'' \\
$[419393]$
\end{quote}

In some cases, the firm went beyond noting a decline in supply or capacity to express more reductions are needed.

\begin{quote}
``We currently anticipate these conditions will require the industry to rationalize production to stabilize the markets so fundamentals can improve. While some supply response has already occurred, more reductions are necessary to balance the market.'' \\
$[185218]$
\end{quote}

\begin{quote}
``On the market side, [there is] still [a] challenging environment in Europe. Industry-wide capacity closures are supportive [and] if we can see maybe 1 or 2 more machine closures in Europe, I think then we are more or less moving it.'' \\
$[260640]$
\end{quote}

\begin{quote}
``There continues to be too much capacity in the global preclinical services industry, and the excess capacity keeps pressure on pricing. We recognized that we would have to reduce capacity in order to improve utilization. A competitor has just announced capacity reductions, and we expect that the pharma industry will continue its efforts to sell and/or close preclinical space. We believe that more capacity remains to be rationalized [and] these actions will help to improve the pricing dynamics.'' \\
$[389910]$
\end{quote}

In nine of the 40 transcripts, reduced industry supply or capacity is mentioned along with references to competitors doing so and the firm supporting such conduct.

\begin{quote}
``We are limiting our investments with reductions in capital [and] production declines will soon be evident. We are witnessing peer companies follow our lead [and] think more companies should adopt this prudent practice.'' \\
$[419394]$
\end{quote}

\begin{quote}
``Industry linerboard and Kraft paper capacity will be substantially reduced [and] we should expect to see further capacity closures announced in the coming months. [This] should lead to substantial price restoration [as] high operating rates give the industry courage to raise prices. I believe that's good for our shareholders and good for the industry.'' \\
$[597969]$
\end{quote}

\begin{quote}
``The industry is oversupplied in NAND, and we need to slow down bit supply. So everybody is slowing. I think the industry is reacting well. We're not trying to gain share. It's about margin profile and profitability.'' \\
$[416422]$
\end{quote}

And when competitors have not reduced supply or capacity, the firm calls them out for not doing so.

\begin{quote}
``We've also done what we think is the prudent thing to do [which] is to effectively provide ... incremental storage to the market by curtailing our production and not pulling it out of the ground. It has been a mystery to me why more producers have not made that same economic choice. We think that it has a chance to restore gas prices, but we'll see how that pans out.'' \\
$[381391]$
\end{quote}

In some instances, a firm notes that it is acting as a market leader in reducing supply or capacity. With that leadership reference comes the implicit invitation for competitors to be followers.

\begin{quote}
``We are currently idling temporarily ... approximately 40\% [of] operating capacity [and] we will be implementing a price increase of up to 10\%. There is a supply and demand imbalance in this business, and being a market leader, we have taken a step to reduce our capacity.'' \\
$[261028]$
\end{quote}

\begin{quote}
``We took our production level down [and] we believe our actions had a positive effect on bringing industry ethanol stocks into a more stable supply and demand balance. Somebody had to exert leadership in this industry to do it. And if we did not slow down, we would be in a much worse margin situation today.'' \\
$[348864]$
\end{quote}

\begin{quote}
``We have made production curtailments. Following our lead, we're now seeing a number of recent capacity curtailments by other producers as the industry takes action necessary to support a recovery in prices. We view this trend favorably and see this dynamic as critical to stabilizing and eventually improving the pricing environment.'' \\
$[569063]$
\end{quote}

\subsubsection{Category 2: The industry needs to reduce supply or capacity or to show discipline}
The second most common category is when a firm announces that the industry needs to reduce supply or capacity or to show discipline. These instances made up 28\% of the true positives. As documented in airlines and steel (see Aryal et al (2022) and Harrington (2022)), ``discipline'' is a common code word for ``less competition''. Other loaded words include firms being ``responsible'' and ``rationalizing'' capacity.

\begin{quote}
``We do consider the possibility of reducing our production in 2017 in order to have a better market environment and also a better environment for pricing. If we in the industry … do not have a discipline regarding our supply, ... we could be generating a situation in which the industry's ROIC could be lower. We are not sure that this will bring other companies to follow in terms of the supply strategies. But what is clear to us is that [we have] the responsibility of showing the facts to be followed in order to seek profitability for the whole industry.'' \\
$[406151]$
\end{quote}

\begin{quote}
``It is important that the U.S. oil and gas industry as well as other global producers continue exerting capital discipline to not overproduce. Everyone needs to continue to cutback and participate as this slump continues. We reduced our pace of growth to deliver a disciplined approach to value creation.'' \\
$[325885]$
\end{quote}

\begin{quote}
``We plan to optimize near-term pricing through the voluntary curtailment of approximately 100 million cubic feet per day over the next several months. Our hope is that our peers, some by choice rather than necessity, will do the same. In today's market, I believe what our industry needs is a more measured pace of growth.'' \\
$[228511]$
\end{quote}

One firm went so far as to specify how much capacity needed to be cut.

\begin{quote}
``It's an obvious conclusion that scrapping has to take place [and] according to my calculations, about 1/3 of the global fleet should be scrapped. Consolidation is inevitable and will take place in our industry.'' \\
$[494779]$
\end{quote}

These expressions of an industry need to cut supply or capacity may be explained in the context of raising prices.

\begin{quote}
``The average prices over the last 15 years in Europe has been EUR 550 … and we believe that ... the trend prices should be more in this range than the historically very low levels that we've seen this year of EUR 420, EUR 430.'' \\
$[260628]$
\end{quote}

The firm may go on to note being encouraged by competitors’ conduct.

\begin{quote}
``The market has overcapacity now, which puts the leverage in the customers' hands. You can resolve overcapacity by not making investments. What needs to happen is capacity needs to shrink. I've been encouraged by what I see our competitors doing.'' \\
$[199854]$
\end{quote}

Especially egregious is when a firm announces how its own supply or capacity cuts reflect ``doing its part by exerting discipline'' [352688] and then demands competitors to pull their weight.

\begin{quote}
``We've implemented planned supply discipline. We think we've done our share and maybe more.'' \\
$[519067]$
\end{quote}

\begin{quote}
``And candidly, we can't do it alone. It's got to be an industry level. There's a gross oversupply in the space. We're taking aggressive action on the supply side.'' \\
$[508184]$
\end{quote}

\begin{quote}
``The market is clearly oversupplied. We are reducing capital by nearly 20\% and production approximately 15\%. One company is not positioned to `fix the market'.'' \\
$[381486]$
\end{quote}

\begin{quote}
``We proactively impaired and early retired a number of our older coastal barges and tugboats. We believe this action is necessary not only to help restore an early balance into the market, but also to improve our profitability going forward. If others will follow our lead, an accelerated balance can be restored to the struggling market. We did our part to rationalize the industry's fleet and improve our profitability going forward, but we need the industry to mirror our actions.'' \\
$[327069]$
\end{quote}

\subsubsection{Category 3: The industry has shown or will show supply or capacity or price discipline}
Rather than expressing a need for more supply or capacity discipline, a firm may note that it has occurred or will occur and expresses support for it. This or similar language appeared in 11\% of the true positive transcripts. These transcripts note the firm’s supply or capacity cuts in the context of an industry effort.

\begin{quote}
``Market conditions remain challenging, exacerbated by the overall oversupply in the industry. We started maintenance of our facilities and adjusted our production utilization rate to 50\% in light of weak market demand. There has been consensus among the manufacturers to promote a healthier development of the industry through exercising self-discipline.'' \\
$[396666]$
\end{quote}

\begin{quote}
``We have seen proactive initiatives to restore the industry's healthy development. Implementing sales discipline would be fundamental to mitigating the irrational competition amid falling prices. We plan to maintain a relatively low utilization rate in 2025 until a turning point emerges in the sector.'' \\
$[396668]$
\end{quote}

\begin{quote}
``If we could scare people out, and they knew how much we were going to be building [then] maybe they wouldn't build theirs. We were trying to get some attention and show that, hey, everybody else should be disciplined, and we're going to come out and throttle this. This is not a cycle where you would expect to see a tremendous amount of capacity come on.'' \\
$[196851]$
\end{quote}

A firm may commend competitors for acting responsibly with respect to their output or prices.

\begin{quote}
``We finally have sanity back in the seaborne iron ore market. I truly commend Rio Tinto and Vale for eliminating the reckless behavior that had infected the market for a number of years. We expect this type of responsible behavior to carry on, which will continue to support strong iron ore prices.'' \\
$[445959]$
\end{quote}

\begin{quote}
``The entire industry moved to action and acted responsibly last year. So we were able to, for the first time, put the prices up and that's a great thing for our business and for our industry.'' \\
$[558080]$
\end{quote}

Based on a recent change in market structure, a firm may forecast future discipline.

\begin{quote}
``We feel that consolidation somehow will improve market discipline going forward. Given the fact that operating in such a competitive scene has led to price erosion in data service pricing, improving market discipline is a very welcome element and consequence of this transaction.'' \\
$[309604]$
\end{quote}

\subsubsection{Category 4: The industry needs to raise prices}
Rather than communicating that the industry needs lower supply or capacity, a firm may announce that it needs higher prices. Such was present in 11\% of the true positive transcripts.

\begin{quote}
``Tariff repair is sorely needed for return ratios to improve. We could do it, we could lead it, we could start it, but the fact is that if competition doesn't follow, then it will hurt us. All the players need it.'' \\
$[250287]$
\end{quote}

\begin{quote}
``I think it's not only our desire to generate that level of profitability. So the need of increasing pricing, I think, will be more an industry need. We think we've arrived at a time where it was necessary to increase a little bit our pricing.'' \\
$[408614]$
\end{quote}

\begin{quote}
``We have an intention to do further price increases. The entire industry has the same need. Otherwise it wouldn't work. So there is a strong need for price increases in Consumer Tissue in Europe. Many of our competitors in Europe have lower margins than we have to start with. So they are in as much need of price increase as we are or even more so.'' \\
$[215003]$
\end{quote}

The called-for price increases may be in the context of responding to cost increases or ending a price war.

\begin{quote}
``We as an industry need to do a better job of tying those contracts to a CPI that actually means something to us, right? We can't be the volume leader by stealing from our competition because we know what that's going to lead to. But in a market that's growing, we should get our fair share of volumes.'' \\
$[375388]$
\end{quote}

\begin{quote}
``Price war destroys value, and it should be reversed soon. If you look at the structural positioning of Telecom Italia, we are by far the strongest player. And therefore I think that the others are suffering as much and more than we are suffering. So it is not something that can be sustainable. This price war is unsustainable for all the players. I consider it a missing opportunity for the whole Italian market.'' \\
$[305524]$
\end{quote}

\subsubsection{Category 5: The firm is being or will be less aggressive in pricing and references competitors’ conduct in a manner encouraging similar conduct}
In eight transcripts, the firm announced it has raised price or will raise price and, when referencing competitors’ pricing, encouraged them to do so or commended them for having done so.

\begin{quote}
``Our strategy of increasing the prices during quarter one has paid off. We [are] going to increase prices further in second half [as] we believe that there is still room to go in that direction. We hope that the market goes in the same direction. The market is now rational, and we believe that there is room for it to keep improving.'' \\
$[309638]$
\end{quote}

\begin{quote}
``We don't want to adjust prices down. We are not cutting prices further. We have been able to maintain our prices. We have seen signs that some of our competitors are doing the same.'' \\
$[223177]$
\end{quote}

\begin{quote}
``Our primary objective will be optimizing profits. We made a conscious decision to trade off attendance [i.e., volume] for yield. We can lean even more heavily into pricing and expect this will further lower our attendance. I can congratulate SeaWorld and Cedar Fair for how they have maintained better pricing integrity than us.'' \\
$[526439]$
\end{quote}

The firm may recognize its role as a market leader with regards to pricing.

\begin{quote}
``We adopted a very clear strategy of price-over-volume strategy for the bromine [and] we are acting as a market leader. When there is a need, we are reducing … our volumes in order to keep the prices at the right level. We don't want prices of bromine as the leader in the market to be too high because this will start to create incentives for other producers. We believe that the current price level of the bromine is more or less optimal.'' \\
$[631460]$
\end{quote}

Some other examples include a firm putting forth an industry objective that puts more weight on price and margins than on market share and volume. If adopted by firms, it has the implication of less aggressive competition and higher prices.

\begin{quote}
``Not only Seagate but the entire industry, decide to focus more on improving the bottom line than fighting for market share. Market share is not our focus. It's taking the right profit from this product. Industry could be constrained. Our objective is to be very reasonable with the capacity we put in place.'' \\
$[408622]$
\end{quote}

\begin{quote}
``Value over volume. The industry is not really equipped to really capture that value proposition. We're going to have to start, as an industry … to collaborate on turning this dynamic around [by] consolidating districts, rationalizing capital expenditures, not building duplicate infrastructure in the same camps.'' \\
$[325064]$
\end{quote}

Finally, one transcript facilitated hub-and-spoke collusion as a manufacturer-hub sought to dampen competition among dealer-spokes.

\begin{quote}
``We will not oversupply the market [and] we hope that our dealers will follow suit and stop competing with themselves.'' \\
$[548932]$
\end{quote}

\subsubsection*{Summary}
The pool of true positives provides encouraging evidence for an LLM being part of a screening protocol. Out of hundreds of thousands of transcripts, the LLM managed to identify a small number that contained collusive content and would be worthy of further investigation. In addition, it did so for content spanning many different forms – reducing supply, constraining capacity, raising prices, stemming price decreases – with messages that conveyed what the industry needs to do, what it did and should continue to do, and a firm acting a leader for the industry. At the same time, it has its limitations as reflected in many false positives. We next turn to examining those false positive transcripts towards understanding how the LLM was misled.

\subsection{How the LLM Got It Wrong}
Out of the 301 LLM-validated transcripts with the highest scores, 192 were classified by the human auditor as false positives. Of those, 29 of them lacked any relevant content and have been interpreted as hallucinations. The other 163 contain some content which, in principle, could be part of a collusive message. Specifically, the transcripts refer to the firm’s competitive conduct or the state of market competition though did not connect it to the conduct of competitors or the industry and thereby fell short of facilitating coordinated conduct. Here we describe the most common content that appeared in the LLM’s excerpts and appeared to be the justification for the LLM unjustifiably assigning it a high score.

67 transcripts referred to a firm reducing, restricting, or managing its supply or capacity in a manner consistent with higher prices or avoiding lower prices. However, it did not say or suggest that this conduct is something the industry should or will or needs to do.

\begin{quote}
``Our intention [is] to maintain a disciplined sales approach. Rather than selling at material decrease prices. We implemented a strategy to hold back product [and] resulting in a significant buildup of the inventories.'' \\
$[298141]$
\end{quote}

\begin{quote}
``We will continue to manage our customer needs and will not restart capacity on a whim just to add tonnage to the spot market. That would not be good for anyone. Value over volume is a simple philosophy that we will carry on into the future, and it's why you are not going to hear me talk about capacity utilization or even market share.'' \\
$[445982]$
\end{quote}

\begin{quote}
``We stand by our strategy of not forsaking value for volume and protecting our profitability with continued actions on our cost structure [and we] cut our production capacity.'' \\
$[458895]$
\end{quote}

\begin{quote}
``We go down by elevator and we go up by stairs. So to avoid this huge volatility, we are aiming … to reduce capacity.'' \\
$[198201]$
\end{quote}

By noting its leadership position or forecasting constrained industry supply, some announcements do get close to suggesting that firms coordinate on reducing supply or capacity.

\begin{quote}
``Our first step is to stop the significant price declines. The most important tool that we have to accomplish this is to reduce capacity to get glass supply in line with market demand. When you're more than half the market, that really adds up.'' \\
$[252936]$
\end{quote}

\begin{quote}
``And until now, we have not heard anyone else announcing any capacity enhancements in the short term. None of them is increasing capacity.'' \\
$[452530]$
\end{quote}

\begin{quote}
``Additional capacity reductions can be expected [and] we forecast an extended period when the … industry operates a smaller, more profitable footprint. Our outlook is for lumber pricing to gradually increase.'' \\
$[654876]$
\end{quote}

48 of the transcripts refer to higher prices - whether noting having raised prices or expressing a need to raise prices - or to pricing power. Such a statement might be made in the context of reducing supply.

\begin{quote}
``We reduced our … production capacity by more than 20\% [and] we have already announced another 20\% price hike. We have started to see the sign of improving pricing trend.'' \\
$[638665]$
\end{quote}

\begin{quote}
``We're going to... be price stewards. We don't chase volume with price.'' \\
$[422806]$
\end{quote}

\begin{quote}
``We have taken a call … to increase prices of cement to help manage the rising cost of production.'' \\
$[320434]$
\end{quote}

With a reference to competitors reducing capacity in the context of an emphasis on raising prices, this announcement is getting close to having collusive content.

\begin{quote}
``In the short term, there's only one priority, and that's price increases and price increase and price increases. We have good capacity utilization. We have strong brands. We always go for pricing every time. Many competitors, especially in the tissue area, announced that they were shutting down paper machines due to this is looked at an Italian player ... and I think we'll see more of that.'' \\
$[215021]$
\end{quote}

In several transcripts, a firm announced it was striving to avoid lowering prices.

\begin{quote}
``We concluded that lowering price to gain share in a market like we are experiencing this year was not the right strategy.'' \\
$[375619]$
\end{quote}

\begin{quote}
``We don't want to underprice and overshoot on demand. We keep some of the exclusive capacity in order to be able to sell at high prices. We want to be profitable [even] if we lost a couple of market share percentage points.'' \\
$[257163]$
\end{quote}

\begin{quote}
``[We are] idling significant chunks of our assets and slashing our participation in weak markets. We're not just going to participate in those weak markets because it's a chase down into the mud.'' \\
$[662160]$
\end{quote}

A firm indicated a strategy to strengthen pricing power.

\begin{quote}
``We deliberately allocate inventory to fiscal years and regions to ensure we maintain apparent scarcity and, therefore, protect our brand health and protect our pricing power.'' \\
$[631021]$
\end{quote}

\begin{quote}
``We're in the early stages of having a terrific opportunity over the next 5 years to be in a position where we can increase the returns at our disposal facilities, primarily through price. We’re one of the only players with a little bit of excess capacity left.'' \\
$[190678]$
\end{quote}

A firm described how it is being more disciplined and less aggressive which, if it had expressed competitors acting or needing to act in a similar manner, would have encompassed collusive content.

\begin{quote}
``We do not want to accelerate a rapid decline in price of the products. We will only ramp it up very moderately when we start production because we do not want to accelerate a collapse in price. We want to be more disciplined in our approach.'' \\
$[604617]$
\end{quote}

\begin{quote}
``We are not chasing market share. We are focusing on quality customers and then driving spend on our network by our existing customers.'' \\
$[643679]$
\end{quote}

\begin{quote}
``We have … implemented Microchip's disciplined pricing process … to ensure that we are competitive and don't chase bad business in the pursuit of profitless prosperity.'' \\
$[620228]$
\end{quote}

In this case, the firm is expressing an aspiration for the industry which is getting close to collusive content.

\begin{quote}
``We would like to drive the industry in a direction where we can move back to more positive margins.'' \\
$[550676]$
\end{quote}

A firm noted that it has market power or what the sources of market power are. While in one case it does mention reduced industry capacity being a source, it is of the form of a factual observation.

\begin{quote}
``Supply is having a harder time to react to demand [and] that's [where] we're getting the pricing power.'' \\
$[401404]$
\end{quote}

\begin{quote}
``When our actual sales pace exceeds our budgeted sales pace, we try to push prices.'' \\
$[651762]$
\end{quote}

\begin{quote}
``We expect to continue to implement favorable pricing actions … due to capacity constraints in our industry. We are operating today in a very strong demand environment without enough industry capacity to meet that demand and expect that to continue into the future.'' \\
$[593253]$
\end{quote}

There are some transcripts referring to a reduced industry supply or capacity but only in the context of stating it as a fact. Still, the announcements are getting close to content that may encourage competitors to reduce supply or capacity.

\begin{quote}
``This moment of rationalization of supply is new in the history of the Brazilian civil aviation industry. We wouldn't like to start or restart price war in the Brazilian market. Our strategy of reducing capacity is indeed in place and it's showing results.'' \\
$[481706]$
\end{quote}

\begin{quote}
``[There is] limited new barge construction in the industry and many units going in for maintenance. So we should have a net decline this year in supply. Nobody is really panicking and trying to go out and build any new equipment right now.'' \\
$[327107]$
\end{quote}

In 14 transcripts, the firm refers to industry consolidation by noting that it has been consolidating, needs to consolidate, or will consolidate. In some instances, the firm attributes consolidation to too much supply or capacity or too many firms. In one case, they do seem to be facilitating such consolidation.

\begin{quote}
``We expect the industry will consolidate to approximately 5 Pan-India players. Once that happens, the industry will further see tariff hardening, with pricing power returning.'' \\
$[321900]$
\end{quote}

\begin{quote}
``In order to really improve value creation, we need consolidation and further restructuring. There’s a need for further consolidation and restructuring in this industry to be able to really actually take the costs down and balance the markets as well.'' \\
$[414209]$
\end{quote}

In summing up, the LLM identified instances in which a firm’s message was consistent with less competition in that it referred to reducing supply or capacity or raising or at least not lowering prices. To facilitate collusive conduct, a firm needs to link such content to industry conduct in a manner that encourage similar conduct by competitors. The LLM failed in not verifying that such a link was present in the firm’s message. Consequently, there was a high false positive rate in the transcripts given the highest scores by the LLM.

\subsection{Takeaways}
Out of hundreds of thousands of transcripts, the LLM was able to pull out a select few containing collusive content. This content takes a variety of forms including expressing a need for reduced industry capacity, noting that it has taken the lead in cutting supply and saying competitors need to do their part, and putting forth a call for the industry to show pricing discipline. At the same, there is a high rate of false positives. The LLM often did not catch the nuanced difference between a firm describing its own supply or capacity reduction, showing pricing discipline, or some other form of constrained conduct and the firm calling on the industry to do so.

While more exploratory work is needed, the evidence presented here supports a valuable role for an LLM in a screening process. An LLM can substitute for human auditors in performing an initial sifting through of a large corpus of transcripts. However, given the high rate of false positives and the critical importance of a competition agency avoiding the allocation of resources to non-collusive cases, human oversight of the LLM’s output is critical.

Here is a possible protocol:
\begin{description}
    \item[Step 1:] LLM reads and rates a corpus of transcripts.
    \item[Step 2:] For those transcripts with the highest scores, human auditors read the LLM-generated excerpts and identifies those with collusive content.
    \item[Step 3:] For those identified as having collusive conduct by the human audit, human auditors read the entire transcript and determines if it is worthy of further examination.
    \item[Step 4:] For those deemed worthy of further examination, all transcripts are collected for the company along with competitors’ transcripts and read by human auditors.
\end{description}

Based on the assessment of that body of transcripts, it can be determined if an official investigation should be initiated. As the LLM is improved in its performance, it will play more of a role by reducing the effort required in step 2 and assisting with step 4.

\section*{Appendix A: Companies with True Positive Transcripts (Number = 71)}
\begin{itemize}
    \item Air France-KLM (3 transcripts)
    \item AirTran
    \item Alliance Resource Partners
    \item American Airlines (2)
    \item Ampco-Pittsburgh
    \item Atwood Oceanics
    \item Avis Budget
    \item Bharti Airtel
    \item Boral
    \item Cameco
    \item Casella Waste
    \item Cementir
    \item CEMEX
    \item Charles River Laboratories
    \item Cleveland-Cliffs (3)
    \item CNX Gas
    \item Coca-Cola Bottlers Japan (4)
    \item Companhia Sider
    \item Continental Resources
    \item Corning
    \item Daqo New Energy (2)
    \item Delta Airlines (2)
    \item Embraer
    \item Essity
    \item Expand Energy (2)
    \item Ferroglobe
    \item Frontier (2)
    \item Glatfelter
    \item Glencore (2)
    \item Goldcorp
    \item Gol Linhas (7)
    \item Green Plains
    \item Gulfport Energy
    \item Harley-Davidson
    \item Huntsman
    \item ICL
    \item Italcementi
    \item JinkoSolar
    \item Kirby
    \item Micron Technology (6)
    \item Norske Skogindustrier (4)
    \item O-I Glass
    \item Oladapco
    \item Ovintiv
    \item PGS
    \item Pilgrim’s Pride
    \item ProFrac Holding
    \item Prosafe
    \item PT XLSMART Telecom (2)
    \item Reliance Communications (2)
    \item Sappi
    \item Seadrill
    \item Seagate Technology (2)
    \item Six Flags Entertainment
    \item Smithfields Foods (3)
    \item STEP Energy
    \item Superior Energy (2)
    \item Suzano
    \item Telecom Italia
    \item Teva Pharmaceuticals
    \item Tidewater
    \item Trip.com Group
    \item Ultra Petroleum (2)
    \item United Airlines (2)
    \item US Airways (3)
    \item Vodafone
    \item Waste Management
    \item Western Digital
    \item WestRock Paper and Packaging (2)
    \item WH Group
    \item Wienerberger
\end{itemize}

\section*{Appendix B: Categorization of Collusive Content}

\textbf{Industry reduced supply or capacity and references competitors’ conduct or the industry’s interests in a manner supportive of such conduct (40 transcripts)} \\
508031, 597966, 505840, 508027, 419393, 352683, 508036, 352685, 392548, 481676, 481654, 602956, 419394, 481658, 481659, 174683, 416422, 524742, 392547, 481690, 427773, 591146, 597969, 259613, 363812, 381391, 185218, 389910, 634653, 260640, 346016, 611671, 569063, 261028, 348864, 340019, 508199, 215550, 190662, 445956

\textbf{Industry needs to reduce supply or capacity or show discipline (30 transcripts)} \\
183145, 418022, 267433, 244726, 363814, 519067, 174682, 427772, 627577, 406151, 481652, 325885, 260639, 611639, 494779, 508184, 381486, 617313, 352688, 342599, 199854, 227682, 399246, 260628, 199856, 436918, 260643, 228511, 327069, 617015

\textbf{The industry has shown or will show supply or capacity or price discipline (12 transcripts)} \\
196851, 392562, 227685, 399267, 396666, 445959, 508185, 264077, 309604, 558080, 396668, 431603

\textbf{Industry needs to raise prices (12 transcripts)} \\
250287, 408614, 215003, 348098, 558083, 558091, 321896, 300093, 375388, 363809, 305524, 324730

\textbf{The firm is being or will be less aggressive in pricing and references competitors’ conduct in a manner encouraging such conduct (8 transcripts)} \\
558086, 377924, 375653, 309638, 223177, 526439, 631460, 321901

\textbf{Miscellaneous (7 transcripts)} \\
States an industry objective that will tend to lead to higher prices [325064, 408622] \\
Industry will limit capacity or supply or not overproduce [252937, 445983] \\
All firms raised prices and pricing strength is expected to continue [257344] \\
Firm welcomes consolidation, whether it's done by us or others as it will add to discipline and make a better market for all of us [360150] \\
Firm will not oversupply and hopes its dealers stop competing with themselves [548932]